\section{Перспективы дальнейшей работы}
\label{sec:Chapter8} \index{Chapter8}

Развитие темы возможно по нескольким направлениям. Первое~--- тесная привязка модели ресурсов к конкретной микроархитектуре целевого чипа (точные квоты PHV, ширины кроссбаров, ограничения параллелизма action-блоков) и сравнение результатов прототипа с промышленным компилятором на реальных P4-программах. Второе~--- включение более богатых эвристик и постоптимизаций после жадной фазы: локальный перенос таблиц между стадиями, симулированный отжиг на небольшом окне узлов TDG или использование ограниченного целочисленного программирования на кластерах связных таблиц. Третье~--- формализация и доказательство аппроксимационных свойств для упрощённых подклассов графов, что может обосновать выбор порядка обхода в FFLS-подобных схемах.
